\chapter{Tests, déploiement et bilan du stage}
\label{ch:tests-deploiement-bilan}

\section*{Introduction du chapitre}
Ce dernier chapitre décrit la stratégie de validation de l'application,
les modalités de son déploiement, puis dresse le bilan personnel et
professionnel du stage.

\section{Stratégie de tests}

La qualité du backend est vérifiée par une suite de tests unitaires et de
service exécutée avec JUnit~5 et Mockito. Aucun contexte Spring complet
n'est démarré pour ces tests : les dépendances de chaque service sont
simulées avec \texttt{@Mock} ou \texttt{Mockito.spy()}, ce qui rend la
suite rapide à exécuter aussi bien en local qu'en intégration continue, et
permet de tester en isolation la logique métier sans dépendre d'une base
de données réelle. Ce choix a en outre orienté certains choix
d'implémentation présentés au chapitre~\ref{ch:realisation} : les méthodes
\texttt{callGemini} et \texttt{callGrok} de
\texttt{QuestionGenerationService}, par exemple, sont délibérément
déclarées \emph{package-private} plutôt que \texttt{private}, dans le seul
but de pouvoir être simulées par \texttt{Mockito.spy()} sans effectuer de
véritable appel réseau lors des tests.

\subsection{Couverture côté backend}

\begin{table}[H]
  \centering
  \caption{Classes de test backend (\texttt{backend/src/test/java/ma/iatacademy/api/})}
  \label{tab:tests-backend}
  \begin{tabular}{@{}ll@{}}
    \toprule
    Classe de test & Objet couvert \\
    \midrule
    \texttt{QuizScoringLogicTest} & Logique de correction automatique des quiz \\
    \texttt{SecurityRbacTest} & Contrôle d'accès par rôle (RBAC) \\
    \texttt{AgendaServiceTest} & Planification / agenda \\
    \texttt{AssignmentServiceTest} & Devoirs \\
    \texttt{AuthServiceTest} & Authentification, inscription, mot de passe \\
    \texttt{BadgeServiceTest} & Gamification (badges) \\
    \texttt{CertificateServiceTest} & Délivrance et vérification de certificats \\
    \texttt{CommentServiceTest} & Commentaires de leçon \\
    \texttt{CourseAssistantServiceTest} & Assistant de cours (IA) \\
    \texttt{EarlyWarningServiceTest} & Détection d'inactivité \\
    \texttt{GradebookServiceTest} & Carnet de notes \\
    \texttt{GroupServiceTest} & Groupes d'apprenants \\
    \texttt{JobOfferServiceTest} & Bourse à l'emploi (accès diplômés) \\
    \texttt{LessonLockServiceTest} & Verrou d'édition collaboratif \\
    \texttt{LessonNoteServiceTest} & Notes personnelles de leçon \\
    \texttt{MediaServiceTest} & Gestion des médias / extraction PDF \\
    \texttt{MessagingServiceTest} & Messagerie interne \\
    \texttt{ProgressionServiceTest} & Déblocage séquentiel du parcours \\
    \texttt{QuestionGenerationServiceTest} & Génération de questions par IA \\
    \texttt{RateLimitServiceTest} & Limitation de débit Redis \\
    \texttt{SearchServiceTest} & Recherche globale \\
    \texttt{StageServiceTest} & Gestion du stage \\
    \texttt{UfValidationServiceTest} & Validation d'unité de formation \\
    \bottomrule
  \end{tabular}
\end{table}
\tabdesc{Au total, cette suite comporte 138 méthodes de test
(\texttt{@Test}), couvrant en priorité les services identifiés au
chapitre~\ref{ch:realisation} comme porteurs de la logique métier la plus
sensible : notation, sécurité, validation d'UF, intégration à des services
d'intelligence artificielle externes.}

Cette priorisation n'est pas anodine : un service comme
\texttt{ProgressionServiceTest}, par exemple, teste explicitement les deux
branches distinctes du calcul de progression --- apprenant individuel
suivant la règle séquentielle automatique par UF, et apprenant rattaché à
un groupe suivant les assignations manuelles du Directeur --- puisqu'une
régression sur l'une de ces deux branches serait particulièrement
difficile à détecter par simple observation manuelle de l'application, la
seconde branche ne concernant qu'une catégorie précise d'apprenants.

\subsection{Couverture côté frontend et tests de charge}

\begin{table}[H]
  \centering
  \caption{Tests frontend}
  \label{tab:tests-frontend}
  \begin{tabular}{@{}lll@{}}
    \toprule
    Type & Outil & Fichiers \\
    \midrule
    Tests unitaires / composants & Vitest + React Testing Library &
      \texttt{QuizTimer.test.tsx}, \texttt{QuizAttemptReview.test.tsx},
      \texttt{Loader.test.tsx}, \texttt{programme.test.ts},
      \texttt{toast-store.test.ts} \\
    Tests de bout en bout & Playwright & \texttt{auth.spec.ts}, \texttt{quiz.spec.ts} \\
    \bottomrule
  \end{tabular}
\end{table}

Le projet comprend enfin un dossier \texttt{loadtest/} dédié aux tests de
montée en charge, écrits comme des scripts autonomes plutôt qu'intégrés à
la suite JUnit, dans la mesure où leur objet n'est pas de vérifier une
règle métier mais un comportement sous concurrence. Le script
\texttt{quiz-start-concurrency.js} vérifie précisément le point discuté au
chapitre~\ref{ch:realisation} (\S\ref{sec:correction-hybride}) : que le
verrou consultatif Postgres (\texttt{pg\_advisory\_xact\_lock}) empêche
effectivement la création de deux tentatives de quiz simultanées pour un
même couple utilisateur/quiz lorsque plusieurs requêtes de démarrage
arrivent en parallèle. Le script \texttt{login-rate-limit.js} vérifie de
son côté que le mécanisme de limitation de débit décrit au
\S\ref{sec:rate-limiting} rejette effectivement les tentatives de
connexion au-delà du seuil configuré, et \texttt{read-throughput.js} mesure
le débit de lecture soutenu par l'API sur les endpoints les plus
consultés. Ces scripts, bien que ne remplaçant pas une campagne de tests
de charge outillée (par exemple avec k6 ou Gatling), ont permis de valider
empiriquement --- et non uniquement par lecture du code --- que les
mécanismes de concurrence et de limitation de débit se comportaient
comme attendu sous une charge simultanée réaliste.

\section{Déploiement}

À l'heure actuelle, le fichier \texttt{docker-compose.yml} placé à
la racine du dépôt active uniquement les services PostgreSQL 16 et Redis~7,
respectivement exposés sur les ports hôtes 5433 et 6379 et adossés chacun
à un volume nommé (\texttt{postgres\_data}, \texttt{redis\_data}) garantissant
la persistance des données entre redémarrages de conteneur ; chaque service
embarque en outre un test de santé (\texttt{healthcheck}) --- \texttt{pg\_isready}
pour PostgreSQL, \texttt{redis-cli ping} pour Redis --- permettant à
d'éventuels services dépendants d'attendre leur disponibilité réelle avant
de démarrer. Les services applicatifs (API Spring Boot et frontend
Next.js) sont pour leur part déjà décrits, à l'état de commentaire dans ce
même fichier, en attente d'une activation ultérieure une fois la
configuration de production stabilisée.

Cette topologie cible n'est cependant pas une simple hypothèse de travail :
elle est déjà entièrement anticipée par la configuration nginx du projet
(\texttt{deploy/nginx/conf.d/elearning.conf}), qui définit deux
\emph{upstreams} nommés --- \texttt{iat\_api} pointant vers le conteneur API
sur le port 8080, et \texttt{iat\_frontend} pointant vers le conteneur
frontend sur le port 3000 --- et route les requêtes entrantes en
conséquence : les chemins \texttt{/api/}, \texttt{/swagger-ui/} et
\texttt{/api-docs} sont dirigés vers l'API (avec un délai d'attente étendu
à 300 secondes pour tolérer les opérations longues), tandis que toute autre
requête, y compris les mises à niveau WebSocket, est dirigée vers le
frontend. Cette configuration fixe également une taille maximale de corps
de requête de 520~Mo, dimensionnée pour absorber les imports de médias
volumineux (vidéos, PDF) traités par la plateforme. Le backend lui-même
est packagé, via \texttt{backend/Dockerfile}, selon un \emph{build} Docker
en deux étapes : une première image de compilation (\texttt{eclipse-temurin:17-jdk})
qui télécharge les dépendances Maven puis compile l'application, et une
seconde image d'exécution, nettement plus légère
(\texttt{eclipse-temurin:17-jre}), qui ne conserve que l'archive
exécutable finale --- une pratique standard qui réduit sensiblement la
taille de l'image livrée en production par rapport à une image de
compilation unique.

\begin{figure}[H]
  \centering
  \includegraphics[width=0.9\textwidth,height=0.85\textheight,keepaspectratio]{diagramme_deploiement.png}
  \caption{Architecture de déploiement}
  \label{fig:deploiement-detail}
\end{figure}

Sur le plan de la configuration, l'ensemble des paramètres sensibles ---
identifiants de connexion à la base de données, clé de signature des
jetons JWT, clés d'API des fournisseurs d'intelligence artificielle
externes --- est injecté exclusivement par variables d'environnement
(fichier \texttt{.env}, jamais versionné dans le dépôt ; voir
\texttt{.env.example} pour la liste complète des clés attendues et leur
rôle), conformément au principe des \emph{Twelve-Factor Apps} qui veut
qu'une même image d'application puisse être déployée sur des
environnements différents sans qu'aucun secret ne soit jamais gravé dans
le code ou dans l'image elle-même. Ce même mécanisme est celui qui permet,
par exemple, d'activer ou de désactiver silencieusement l'intégration à
Gemini et à Grok décrite au \S\ref{sec:generation-ia} : un fournisseur est
considéré comme configuré dès lors que sa clé d'API n'est pas vide,
sans qu'aucun indicateur booléen supplémentaire ne soit nécessaire.

\section{Bilan du stage}
\label{sec:bilan-difficultes}

\subsection{Difficultés rencontrées}

Ce travail n'a pas été mené sans difficultés. La principale d'entre elles
a consisté à traduire des règles métier parfois implicites --- notamment
la distinction entre les UF qui exigent une validation manuelle du
Directeur et celles qui n'en exigent pas, ou encore la coexistence de deux
stratégies de déblocage du contenu selon qu'un apprenant suit son parcours
en ligne de façon autonome ou qu'il appartient à une cohorte encadrée ---
en un modèle de données et une logique applicative à la fois simples et
fidèles à la réalité du terrain. La confrontation régulière avec l'équipe
pédagogique d'\Entreprise{} a été déterminante pour éviter qu'une
simplification excessive du modèle ne finisse par trahir ces règles.

Le chapitre~\ref{ch:realisation} assume également, en toute transparence,
deux limites concrètes identifiées dans le code existant --- une
incohérence dans l'application du contrôle de propriété au niveau du
service de gestion des devoirs, et l'absence d'une reprise automatique
après l'interruption d'une campagne email (\S\ref{sec:limites}) --- qui
constituent naturellement les premières pistes d'amélioration à court
terme plutôt que des défauts dissimulés.

\subsection{Apports techniques}

Ce stage a permis d'approfondir des compétences qui dépassent largement la
simple implémentation de fonctionnalités CRUD : la conception d'un modèle
de données fidèle à des règles métier complexes et partiellement
implicites, l'intégration résiliente de services externes (IA générative,
service SMTP, hébergement vidéo) avec gestion explicite des scénarios
d'échec, la sécurisation d'une API REST à plusieurs niveaux (RBAC et
contrôle de propriété), et la mise en place d'une suite de tests
automatisés couvrant en priorité la logique métier la plus sensible plutôt
que l'ensemble du code de façon uniforme.

\subsection{Apports personnels et professionnels}

Sur le plan personnel, mener seul --- en lien direct et continu avec
l'équipe pédagogique d'\Entreprise{} --- l'intégralité d'un projet logiciel
réel, de l'analyse des besoins jusqu'au déploiement, a été une expérience
professionnelle formatrice à part entière. Elle a permis de mettre en
pratique, sur un projet destiné à des utilisateurs réels, des compétences
d'analyse, de conception et de rigueur technique qui dépassent largement
le cadre d'un exercice académique classique, ainsi qu'une autonomie
accrue dans la priorisation du travail au sein d'une démarche itérative
non encadrée par des cérémonies Scrum formelles.

\section*{Conclusion du chapitre}
Entre une suite de 138 tests backend complétée par des tests frontend et
des scripts de montée en charge ciblés sur les mécanismes de concurrence
les plus sensibles, une architecture de déploiement dont la topologie
cible est déjà anticipée par la configuration nginx du projet, et le
bilan qui vient d'être dressé, ce chapitre boucle la démarche de
validation du projet. Il ne reste qu'à revenir, dans la conclusion
générale, sur l'ensemble du chemin parcouru durant ce stage.
